
%%%%%%%%%%%%%%%%%%%%%%%%%%% ESPECIALES %%%%%%%%%%%%%%%%%%%%%%%%%%%%%%%%%%
% \RequirePackage{expl3}    % \usepackage cannot be used before \documentclass
% \ExplSyntaxOn             % Switch on expl3 syntax
% \pdf_version_gset:n{1.7}  % Use provided expl3 function
% \ExplSyntaxOff            % Switch off expl3 syntax

%%%%%%%%%%%%%%%%%%%%%%%%%%% Tipografía %%%%%%%%%%%%%%%%%%%%%%%%%%%%%%%%%%

\usepackage[sfdefault,t,lf]{carlito}    % Alternativa OpenSource a Calibri, definir sans serif como predeterminado, numerales tabulares
% \usepackage[default,lnum]{opensans}
% \usepackage{CormorantGaramond}

% \usepackage{lucidabr}

% \usepackage{fontspec}

\fontsize{8.8pt}{12pt}\selectfont
\renewcommand{\baselinestretch}{1.2}

% \linespread{1.3}
% \DeclareUnicodeCharacter{2212}{-}

\usepackage{polyglossia}
	\setdefaultlanguage[spanishoperators=all]{spanish}  % Definir la localización del documento como es-MX (para que cuadro sea tabla, etc.)
	% \setdefaultlanguage{mexican}
	\renewcommand{\listtablename}{Índice de tablas}

\usepackage{translations}
	\DeclareLanguage{Spanish}

\usepackage{changelog}  % Control de versiones

\usepackage{pdflscape}
\usepackage{pifont}
\newcommand{\cmark}{\ding{51}}%
\newcommand{\xmark}{\ding{55}}%

% \usepackage[%
% toc=true,
% postpunc={;},
% acronyms,
% xindy]{glossaries-extra}    % Manager de glosarios, usar xindy para el indice alfabético
% % \usepackage{glossary-mcols}
% % \renewcommand{\glsmcols}{2}
% \setglossarystyle{long3col}

% 	\makeglossaries{}
% 	% \setabbreviationstyle[acronym]{short-long}  % Estilo de acrónimos
% 	\setabbreviationstyle[acronym]{long-short}
% 	% \setglossarystyle{mcolindex}
% 	\loadglsentries[main]{../glosario-def.tex}  % Archvio que LuaLaTeX toma para las definiciones

\usepackage[
	record=hybrid,% provides group and location fields (and other stuff)
	nopostdot,
	acronym,
	stylemods=tree]{glossaries-extra}
\setabbreviationstyle{long-short-sc}
\GlsXtrLoadResources[
	src={../glosario.bib},	% bib files
	selection={all},			% select all entries
	sort={es-MX}
]
\makeglossaries


% \setglossarystyle{longragged}


\renewcommand{\glossarypreamble}{\fontsize{10pt}{12pt}\selectfont}

\usepackage{makeidx}
\makeindex

% \usepackage[xindy]{imakeidx}    % Para hacer indices
% 	\makeindex[columns=2]   	% Estilo del indice alfabético
	
\newcounter{note}[section]  % Nuevo contador para notas
\renewcommand{\thenote}{\thesection.\Alph{note}}    % Definir estilo de numeración de las notas
\usepackage[framemethod=TikZ]{mdframed}             % Motor gráfico para notas (TikZ)
\newenvironment{note}[1][]{%                        % Nuevo entorno de Nota (note) 
    \refstepcounter{note}
    \begin{mdframed}[%                                  
        frametitle={Nota \thenote\ #1},                 % Titulo de las notas
        skipabove=\baselineskip{} plus 2pt minus 1pt,     % Espaciado superior
        skipbelow=\baselineskip{} plus 2pt minus 1pt,     % Espaciado inferior
        linewidth=0.5pt,                                % Tamaño de linea del recuadro
        frametitlerule=true,                            % incluir titulo
        frametitlebackgroundcolor=gray!30               % Color del recuadro
    ]%
}{%
    \end{mdframed}
}

\newcounter{contact}[section]                           % Nuevo contador para contactos
\renewcommand{\thecontact}{\thesection.\Alph{contact}}  % Definir que la numeración sea Sección.Alfabético

\newenvironment{contact}[1][]{%
    \refstepcounter{contact}
    \begin{mdframed}[%
        frametitle={Contacto \thecontact\ #1},          % Titulo del contacto
        skipabove=\baselineskip{}plus 2pt minus 1pt,     % Espaciado superior
        skipbelow=\baselineskip{}plus 2pt minus 1pt,     % Espaciado inferior
        linewidth=0.5pt,                                % Tamaño de linea del recuadro
        frametitlerule=true,                            % poner titulo
        frametitlebackgroundcolor=yellow!30             % Color del recuadro
    ]%
}{%
    \end{mdframed}
}

\usepackage[multiple]{footmisc}     % Cambios en la presentación de pies de página
\usepackage{microtype}              % Ajustes de kerning y otros 
\usepackage{graphicx}               % Mejora al motor de imagenes
\usepackage{hyperref}               % Referencias cruzadas en PDF
\usepackage{cleveref}               % Referencias cruzadas en texto
%%%%%% REDEFINICIONES DE REFERENCIAS CRUZADAS %%%%%%%%%
    \crefname{note}{nota}{notas}    
	\Crefname{note}{Nota}{Notas}
	\crefname{contact}{contacto}{contactos}
	\Crefname{contact}{Contacto}{Contactos}
	\crefname{paragraph}{parrafo}{parrafos}
	\Crefname{paragraph}{Parrafo}{Parrafos}
%%%%%%%%%%%%%%%%%%%%%%%%%%%%%%%%%%%%%%%%%%%%%%%%%%%%%%%
\newcommand{\email}[1]{\href{mailto:#1}{#1}}    % Comando para hacer que se puedan clickear emails
\newcommand{\cels}[1]{\qty{#1}{\celsius}}       % Escrir celsius rapido

\usepackage{lastpage}   % Definir variable para el numero total de páginas
\usepackage{siunitx}    % Unidades internacionales y alineamiento de decimales en tablas
	\sisetup{detect-all}    % Detectar automaticamente que la tipografía es sans serif

\usepackage{color}  % definir colores nuevos

\usepackage{tabularray} % Creación de tablas
	\UseTblrLibrary{booktabs,siunitx}   % librerias parea tablas correctas, unidades internacionales
	\definecolor{Gallery}{rgb}{0.929,0.929,0.929}   % definición del color "Gallery"
	\DefTblrTemplate{contfoot-text}{normal}{\textit{Continúa en la siguiente página}}   % Continuación de tabla
	\SetTblrTemplate{contfoot-text}{normal} % Continuación de tabla normal (footer)
	\DefTblrTemplate{conthead-text}{normal}{\textit{(Continuación)}}    % Continuación de tabla
	\SetTblrTemplate{conthead-text}{normal} % Continuación de tabla (header)

    % \usepackage[margin=1cm,bmargin=2cm,tmargin=5cm,headheight=110pt]{geometry}   % Dispoción espacial del documento
% \newgeometry{includeheadfoot, heightrounded, left=1cm, right=1cm, top=1cm, bottom=4cm, headheight=6em, footskip=1.5em}    
%     \savegeometry{L1}

% \newgeometry{includeheadfoot, heightrounded, left=1cm, right=1cm, top=1cm, bottom=4.5cm, headheight=7em, footskip=1.5em, showcrop, showframe}
%     \savegeometry{L2}

\usepackage[%
includeheadfoot,
heightrounded,
left=1cm,
right=1cm,
top=1cm,
bottom=1cm,
headheight=110pt,
footskip=1.5em]{geometry}    

\usepackage[export]{adjustbox}

\usepackage{enumitem}
	\setlist[1]{labelindent=\parindent} % Indentación como parrafo para listas
	\setlist[description]{style=nextline}
	\setlist[itemize]{leftmargin=*} 	% Eliminar margen en izq
	\setlist[itemize,1]{label=---}  	% Viñeta como em-dash
	\setlist[itemize,2]{label=--}
	\setlist{noitemsep} 				% Que no se separen los item
	% [label={\theparagraph--\arabic*}] % Numerarar parrafo-numero

	\newlist{ListaSec}{enumerate}{3}
	\setlist[ListaSec]{label* = \arabic*., name* = \arabic{section}.\arabic{subsection}}


%%%%%%%%%%%%%%%% DEFINICIÓN DE CABECERA %%%%%%%%%%%%%%%%
\usepackage{fancyhdr}   % Usar fancy headers y footers
	\renewcommand{\headrulewidth}{0pt}  % Grosor de regla del header
	\fancyhead[CE,CO,LE,LO,RE,RO]{} % Resetear headers    
    \fancyfoot[CE,CO,LE,LO,RE,RO]{} % Resetear footers

%%%%%%%%%%%%%%%%%%%%%% VARIABLES DE CADA DOCUMENTO DENTRO DE UN PROGRAMA %%%%%%%%%%%%%%%%%%%%%
	\newcommand{\Codigo}{}                  % Codigo del documento
	\newcommand{\FechaPub}{}                % Fecha de publicación
	\newcommand{\Titulo}{}                  % Titulo del documento
	\newcommand{\MayorVer}{}                % Versión mayor
	\newcommand{\MenorVer}{}                % Verisón menor
	\newcommand{\Edit}{\MayorVer.\MenorVer} % Definición de la edición

    \newcommand{\Elaboro}{}                 % Nombre de quien lo elaboró
    \newcommand{\Reviso}{}                  % Nombre de quien revisó el documento
    \newcommand{\Autorizo}{}                % Nombre de quien aprobó el documento

%%%%%%%%%%%%%%%%%% Formato por defecto para informaicón documentada general %%%%%%%%%%%%%%%%%%
\fancypagestyle{formato-PS}{%
	\fancyhead{}
	\fancyfoot{}
	\fancyhead[RE]{%
		\begin{tblr}{%	
			colspec = {X[c,m]X[c,m]},    % Especificación de columnas: Anchura estática, centrado en medio
			hlines,
			vlines,
			width = 0.35\linewidth,}%
			{\textbf{Código:}\\ \Codigo} 				 & {\textbf{Edición:}\\ \MayorVer.\MenorVer} 			\\ 	% Logotipo (2 columnas), red de frios (3 columnas), código del documento
			{\textbf{Fecha de publicación:}\\ \FechaPub} & {Página~\thepage~de~\pageref{LastPage}} 		% Titulo (4 columnas), Fecha, edición
		\end{tblr}
	}%

	\fancyhead[LO]{%
		\begin{tblr}{%	
			colspec = {X[c,m]X[c,m]},   % Especificación de columnas: Anchura estática, centrado en medio
			hlines,
			vlines,
			width = 0.35\linewidth,}%
			{\textbf{Código:}\\ \Codigo} 					& {\textbf{Edición:}\\ \MayorVer.\MenorVer} 			\\	% Logotipo (2 columnas), red de frios (3 columnas), código del documento
			{\textbf{Fecha de publicación:}\\\FechaPub}	& {Página~\thepage~de~\pageref{LastPage}} 		% Titulo (4 columnas), Fecha, edición
		\end{tblr}
	}%

	\fancyfoot[C]{Página~\thepage~de~\pageref{LastPage}}   % Para que muestre Página # de #
	\fancyfoot[L]{\textbf{CONFIDENCIAL}}                    % CONFIDENCIAL
	\fancyfoot[R]{\Codigo~v.\Edit}                          % Estructura del código CODIGO v.EDICION
}%

\fancypagestyle{formato-PI}{%
	\fancyhead{}
	\fancyfoot{}
    \fancyhead[C]{%
	    \begin{tblr}{%	
            colspec = {X[c,m]X[c,m]X[c,m]X[c,m]X[c,m]X[c,m]},    % Especificación de columnas: Anchura estática, centrado en medio
            hlines,
            vlines,
			width = \linewidth,}%
			\SetCell[c=2]{} \includegraphics[height=3em, valign=c]{./RDF_Logo.png}	&	& \SetCell[c=3]{} \large{Red de fríos S.A. de C.V.}	&	&												& {\textbf{Código:}\\ \Codigo}	\\ % Logotipo (2 columnas), red de frios (3 columnas), código del documento
			\SetCell[c=4]{} {\textbf{Titulo:}\\ \Titulo} 							&	&													&	& {\textbf{Fecha de publicación:}\\ \FechaPub}	& {\textbf{Edición:}\\ \Edit}	\\ % Titulo (4 columnas), Fecha, edición
            \end{tblr}
	}
    \fancyfoot[C]{Página~\thepage~de~\pageref{LastPage}}   % Para que muestre Página # de #
    \fancyfoot[L]{\textbf{CONFIDENCIAL}}                    % CONFIDENCIAL
    \fancyfoot[R]{\Codigo~v.\Edit}                          % Estructura del código CODIGO v.EDICION
}%

\setcounter{tocdepth}{4}            % 4 Qué tanto mostrará la tabla de contenidos
\setcounter{secnumdepth}{6}         % Que tantas secciones numerará LaTeX (Se ajustó hasta subparrafo)